\documentclass[11pt]{article}

\usepackage[a4paper,margin=30mm]{geometry}
\usepackage{amsmath,amssymb,amsthm,mathtools}
\usepackage{microtype}
\usepackage[colorlinks=true,linkcolor=blue,citecolor=blue,urlcolor=blue]{hyperref}
\usepackage[nameinlink,noabbrev]{cleveref}

\newtheorem{theorem}{Theorem}[section]
\newtheorem{lemma}[theorem]{Lemma}
\newtheorem{proposition}[theorem]{Proposition}
\newtheorem{corollary}[theorem]{Corollary}
\theoremstyle{definition}
\newtheorem{definition}[theorem]{Definition}
\theoremstyle{remark}
\newtheorem{remark}[theorem]{Remark}

\newcommand{\R}{\mathbb{R}}
\newcommand{\cC}{\mathcal{C}}
\newcommand{\cR}{\mathcal{R}}
\newcommand{\bN}{\boldsymbol{n}}
\newcommand{\bY}{\boldsymbol{y}}
\newcommand{\1}{\mathbf{1}}

\title{Sharp Mesh-Dilation Bounds for Arbitrary Romberg
Extrapolation of Completely Monotone Integrands}
\author{}
\date{}

\begin{document}
\maketitle

\begin{abstract}
Let \(R_{\bN}(f)\) be the value at zero obtained by interpolating
composite trapezoidal approximations \(T_{n_j}(f)\) polynomially in
\(n_j^{-2}\), where \(n_1,\ldots,n_m\) are arbitrary distinct positive
integers. For every nonconstant finite positive Laplace transform \(f\) on a
compact interval and every integer refinement factor \(q\ge2\), we prove the
sharp inequality
\[
q^{-2m}
<
\frac{R_{q\bN}(f)-I(f)}
     {R_{\bN}(f)-I(f)}
<
q^{-1}.
\]
Both constants are approached by exponential integrands, at zero and infinite
decay rate respectively. The ratio bound gives a computable strict enclosure
of the integral between two affine combinations of the coarse and refined
extrapolants. It also yields an exact classification of every affine blend as
a universal upper bound, a universal lower bound, or a rule with no fixed
sign. For twofold refinement, the lower endpoint is midpoint extrapolation and
the upper endpoint gives the sharp midpoint--trapezoid threshold. The proof
uses a positive origin-interpolation formula for complete Bernstein kernels
and a sharp comparison between a decreasing spectral sum and its subsequence
indexed by multiples of \(q\).
\end{abstract}

\noindent\textbf{Keywords.}
Romberg quadrature; Richardson extrapolation; completely monotone function;
Laplace transform; trapezoidal rule; certified integration.

\medskip
\noindent\textbf{MSC 2020.}
65D30; 26A48; 41A05.

\section{Introduction}

Romberg quadrature accelerates the composite trapezoidal rule by polynomial
extrapolation in the squared mesh width. Its classical dyadic table has
positive Peano kernels and supplies one-sided estimates under derivative-sign
hypotheses. Positivity, monotonicity, strict error bounds, and nonclassical
refinement sequences have a substantial literature
\cite{MeirSharma1965,Stroud1965,BulirschStoer1966,
BulirschStoer1967,Lynch1967,Albrecht1972,DahlquistBjorck,Strom1967,
Strom1972,Foerster1982,BrassFischer1999,Fischer2002,BrassPetras}.

This paper considers the cone
\[
f(x)=\int_{[0,\infty)}e^{-tx}\,d\mu(t),
\]
where \(\mu\) is a finite positive measure. By the Bernstein--Widder theorem,
this is the class of completely monotone functions on \([0,\infty)\) with
finite value at zero, restricted to the interval of integration
\cite{Widder,SchillingSongVondracek}. The representation permits a direct
spectral analysis of the complete trapezoidal error rather than an estimate
through finitely many derivatives.

Fix arbitrary distinct panel counts \(n_1,\ldots,n_m\). Let \(R_{\bN}(f)\)
be the origin value of the polynomial interpolating \(T_{n_j}(f)\) in
\(n_j^{-2}\). If every panel count is multiplied by an integer \(q\ge2\), the
error contracts by a factor strictly between \(q^{-2m}\) and \(q^{-1}\).
The lower constant is the local, small-decay scaling of the order-\(2m\)
extrapolation error. The upper constant is the density of multiples of \(q\)
in the spectral lattice. Both constants are sharp.

The ratio theorem gives two computable bounds:
\[
\frac{qR_{q\bN}(f)-R_{\bN}(f)}{q-1}
<
I(f)
<
\frac{q^{2m}R_{q\bN}(f)-R_{\bN}(f)}{q^{2m}-1}.
\]
It also identifies every universally signed affine blend of the two
extrapolants. The upper endpoint is the usual Richardson combination. The
lower endpoint is the extrapolate of the equal-weight rule supported precisely
on the new nodes introduced by \(q\)-fold refinement. For \(q=2\), that rule
is the composite midpoint rule.

The proof has four short components. First, extrapolating
\(y/(y+s)\) from positive nodes to zero produces a positive product. Second,
the exponential trapezoidal error is expanded using \(x\coth x-1\). Third, a
strictly decreasing positive sequence is compared with its \(q\)-multiple
subsequence. Finally, positive Laplace mixtures preserve the resulting ratio
interval.

\section{Quadrature and extrapolation}

Fix \(L>0\). Let \(\cC_L\) be the class of functions
\begin{equation}
f(x)=\int_{[0,\infty)}e^{-tx}\,d\mu(t),
\qquad 0\le x\le L,
\label{eq:laplace-class}
\end{equation}
where \(\mu\) is a finite positive Borel measure. Define
\[
I(f)=\int_0^L f(x)\,dx
\]
and let
\begin{equation}
T_n(f)
=
\frac{L}{n}
\left[
\frac{f(0)+f(L)}{2}
+\sum_{r=1}^{n-1}f\left(\frac{rL}{n}\right)
\right]
\label{eq:trapezoid}
\end{equation}
be the composite trapezoidal rule with \(n\) equal panels.

Let
\[
\bN=(n_1,\ldots,n_m)
\]
consist of distinct positive integers and set \(y_j=n_j^{-2}\). The
extrapolate \(R_{\bN}(f)\) is the value at zero of the unique polynomial of
degree at most \(m-1\) satisfying
\[
P(y_j)=T_{n_j}(f),\qquad j=1,\ldots,m.
\]
No ordering or geometric relation among the \(n_j\) is assumed. For an integer
\(q\ge2\), write \(q\bN=(qn_1,\ldots,qn_m)\).

\begin{remark}
The Lagrange weights at zero are invariant under common positive scaling of
all interpolation nodes. Thus \(R_{q\bN}(f)\) may be evaluated using the
nodes \(q^{-2}y_j\) or, equivalently, the original nodes \(y_j\) with the
trapezoidal error profile dilated by \(q^{-2}\).
\end{remark}

\section{A positive origin-interpolation formula}

For distinct positive nodes \(\bY=(y_1,\ldots,y_m)\), let
\(\Lambda_{\bY}g\) denote the value at zero of the degree-\((m-1)\)
polynomial interpolating \(g\) at those nodes.

\begin{lemma}[Origin interpolation]
\label{lem:origin}
For \(s>0\), let \(h_s(y)=y/(y+s)\). Then
\begin{equation}
\Lambda_{\bY}h_s
=K_{\bY}(s)
:=\prod_{j=1}^{m}\frac{y_j}{y_j+s}.
\label{eq:kernel}
\end{equation}
\end{lemma}

\begin{proof}
Let \(P\) interpolate \(h_s\) at \(\bY\). The numerator in
\[
h_s(y)-P(y)
=\frac{y-(y+s)P(y)}{y+s}
\]
is a polynomial of degree at most \(m\) that vanishes at every \(y_j\).
Consequently,
\[
y-(y+s)P(y)=C\prod_{j=1}^{m}(y-y_j).
\]
Evaluation at \(y=-s\) determines \(C\). Evaluation at \(y=0\) then gives
\[
P(0)=\prod_{j=1}^{m}\frac{y_j}{y_j+s}.
\]
\end{proof}

\begin{proposition}[Complete Bernstein error profile]
\label{prop:cbf}
For \(f\in\cC_L\), define
\begin{equation}
\Phi_f(y)
=
\int_{(0,\infty)}
I_t\left[x_t(y)\coth x_t(y)-1\right]\,d\mu(t),
\label{eq:error-profile}
\end{equation}
where
\[
I_t=\frac{1-e^{-tL}}{t},
\qquad
x_t(y)=\frac{tL\sqrt{y}}{2}.
\]
Then \(\Phi_f\) is a complete Bernstein function with zero constant and
linear terms, and
\begin{equation}
\Phi_f(n^{-2})=T_n(f)-I(f).
\label{eq:profile-samples}
\end{equation}
More precisely,
\begin{equation}
\Phi_f(y)
=
2\int_{(0,\infty)}
I_t\sum_{k=1}^{\infty}
\frac{y}{y+c_tk^2}\,d\mu(t),
\qquad
c_t=\left(\frac{2\pi}{tL}\right)^2.
\label{eq:cbf-representation}
\end{equation}
\end{proposition}

\begin{proof}
For \(f_t(x)=e^{-tx}\), a geometric sum in \eqref{eq:trapezoid} gives
\begin{equation}
\frac{T_n(f_t)}{I_t}
=x\coth x,
\qquad
x=\frac{tL}{2n}.
\label{eq:xcoth}
\end{equation}
The product expansion of \(\sinh x/x\), followed by logarithmic
differentiation, gives
\begin{equation}
x\coth x-1
=2\sum_{k=1}^{\infty}
\frac{x^2}{x^2+\pi^2k^2}.
\label{eq:mittag}
\end{equation}
Substitution yields \eqref{eq:cbf-representation}. This is the standard
integral form of a complete Bernstein function.

All summands are nonnegative, so Tonelli's theorem applies. The bound
\[
0<x\coth x-1<x
\]
also gives
\[
0\le\Phi_f(y)
\le \frac{L\sqrt{y}}{2}\mu((0,\infty)).
\]
Thus \(\Phi_f(y)/y\to0\) at infinity, which rules out a linear term in its
complete Bernstein representation. Equation \eqref{eq:profile-samples}
follows from \eqref{eq:xcoth}, linearity, and exactness on the constant
component corresponding to \(t=0\).
\end{proof}

\begin{corollary}[Positive arbitrary-mesh error]
\label{cor:positive-error}
For \(f\in\cC_L\),
\begin{equation}
R_{\bN}(f)-I(f)
=
2\int_{(0,\infty)}
I_t
\sum_{k=1}^{\infty}K_{\bY}(c_tk^2)
\,d\mu(t).
\label{eq:positive-error}
\end{equation}
In particular, \(R_{\bN}(f)>I(f)\) whenever \(f\) is nonconstant.
\end{corollary}

\begin{proof}
Apply \(\Lambda_{\bY}\) to \eqref{eq:cbf-representation} and use
\cref{lem:origin}. The interpolation functional is a finite linear
combination of finite values, and the resulting representation is positive.
\end{proof}

\section{The sharp mesh-dilation ratio}

The arithmetic estimate below supplies both sharp constants.

\begin{lemma}[Spectral subsequence ratio]
\label{lem:spectral}
Let \(c>0\), let \(q\ge2\) be an integer, and define
\[
a_k=K_{\bY}(ck^2),\qquad k\ge1.
\]
Then
\begin{equation}
q^{-2m}
<
\frac{\sum_{k=1}^{\infty}a_{qk}}
     {\sum_{k=1}^{\infty}a_k}
<
q^{-1}.
\label{eq:spectral-ratio}
\end{equation}
\end{lemma}

\begin{proof}
For every \(k\),
\[
\frac{a_{qk}}{a_k}
=
\prod_{j=1}^{m}
\frac{y_j+ck^2}{y_j+cq^2k^2}
>q^{-2m}.
\]
Summation proves the lower inequality.

The sequence \(a_k\) is strictly decreasing. Partition the positive integers
into consecutive blocks of length \(q\). In the \(k\)-th block,
\[
\sum_{r=1}^{q}a_{q(k-1)+r}>q\,a_{qk}.
\]
Summation over all blocks proves the upper inequality.
\end{proof}

\begin{theorem}[Sharp mesh-dilation bounds]
\label{thm:main}
Let \(f\in\cC_L\) be nonconstant. For every \(m\ge1\), every vector
\(\bN\) of distinct positive panel counts, and every integer \(q\ge2\),
\begin{equation}
\boxed{
q^{-2m}
<
\frac{R_{q\bN}(f)-I(f)}
     {R_{\bN}(f)-I(f)}
<
q^{-1}.}
\label{eq:main}
\end{equation}
Both constants are optimal uniformly over \(\cC_L\).
\end{theorem}

\begin{proof}
For a single exponential \(f_t\), equations
\eqref{eq:positive-error} and the common-scaling invariance give
\[
R_{\bN}(f_t)-I_t
=2I_t\sum_{k=1}^{\infty}K_{\bY}(c_tk^2)
\]
and
\[
R_{q\bN}(f_t)-I_t
=2I_t\sum_{k=1}^{\infty}K_{\bY}(c_tq^2k^2).
\]
Lemma~\ref{lem:spectral} proves \eqref{eq:main} for every \(t>0\).

For a general \(f\), let \(E(t)\) and \(E_q(t)\) denote these two positive
single-exponential errors. Then
\[
\frac{R_{q\bN}(f)-I(f)}{R_{\bN}(f)-I(f)}
=
\frac{\int E_q(t)\,d\mu(t)}{\int E(t)\,d\mu(t)}
\]
is a weighted average of \(E_q(t)/E(t)\), with positive weights
\(E(t)\,d\mu(t)\). Because a nonconstant \(f\) has positive representing mass
on \((0,\infty)\), both inequalities remain strict.

It remains to prove optimality. As \(t\to0^+\), \(c_t\to\infty\), and
\[
c_t^mK_{\bY}(c_tk^2)
\longrightarrow
\left(\prod_{j=1}^{m}y_j\right)k^{-2m}.
\]
Moreover,
\[
c_t^mK_{\bY}(c_tk^2)
\le
\left(\prod_{j=1}^{m}y_j\right)k^{-2m}.
\]
Dominated convergence therefore shows that the ratio for \(f_t\) tends to
\(q^{-2m}\).

As \(t\to\infty\), \(c_t\to0\). Put
\[
g(u)=K_{\bY}(u^2).
\]
The function \(g\) is continuous and integrable on \((0,\infty)\): it is
bounded near zero and \(O(u^{-2m})\) at infinity. Hence
\[
\sqrt{c_t}\sum_{k=1}^{\infty}g(k\sqrt{c_t})
\longrightarrow
\int_0^\infty g(u)\,du,
\]
whereas
\[
\sqrt{c_t}\sum_{k=1}^{\infty}g(qk\sqrt{c_t})
\longrightarrow
\frac1q\int_0^\infty g(u)\,du.
\]
The ratio tends to \(q^{-1}\). Both endpoint constants are therefore sharp
within the exponential subfamily.
\end{proof}

\begin{corollary}[Strict contraction]
\label{cor:contraction}
Under the hypotheses of \cref{thm:main},
\[
I(f)<R_{q\bN}(f)<R_{\bN}(f).
\]
\end{corollary}

\section{Certified bounds and affine blends}

Write
\[
A=R_{\bN}(f),
\qquad
B=R_{q\bN}(f).
\]

\begin{corollary}[Two-sided enclosure]
\label{cor:enclosure}
Every nonconstant \(f\in\cC_L\) satisfies
\begin{equation}
\boxed{
\frac{qB-A}{q-1}
<
I(f)
<
\frac{q^{2m}B-A}{q^{2m}-1}.}
\label{eq:enclosure}
\end{equation}
Both affine coefficients are optimal uniformly over \(\cC_L\).
\end{corollary}

\begin{proof}
Let
\[
r=\frac{B-I(f)}{A-I(f)}.
\]
Solving for the integral gives
\[
I(f)=\frac{B-rA}{1-r}.
\]
Since \(B<A\), the right-hand side is strictly decreasing in \(r\).
Substitution of \(q^{-2m}<r<q^{-1}\) proves
\eqref{eq:enclosure}. Optimality follows from the endpoint families in
\cref{thm:main}.
\end{proof}

For \(\beta\in\R\), define
\[
Q_\beta(f)=(1-\beta)A+\beta B.
\]

\begin{theorem}[Complete affine-blend classification]
\label{thm:blend}
The inequality
\[
Q_\beta(f)\ge I(f)
\]
holds for every \(f\in\cC_L\) if and only if
\begin{equation}
\beta\le\frac{q^{2m}}{q^{2m}-1}.
\label{eq:upper-beta}
\end{equation}
The reverse inequality
\[
Q_\beta(f)\le I(f)
\]
holds for every \(f\in\cC_L\) if and only if
\begin{equation}
\beta\ge\frac{q}{q-1}.
\label{eq:lower-beta}
\end{equation}
For
\[
\frac{q^{2m}}{q^{2m}-1}
<
\beta
<
\frac{q}{q-1},
\]
the error has no universal sign.
\end{theorem}

\begin{proof}
For nonconstant \(f\),
\[
Q_\beta(f)-I(f)
=(A-I(f))\bigl[1-\beta(1-r)\bigr],
\]
where \(r\) ranges sharply over
\((q^{-2m},q^{-1})\). The bracket is nonnegative throughout this interval
exactly under \eqref{eq:upper-beta}, and nonpositive throughout exactly under
\eqref{eq:lower-beta}. If \(\beta\) lies strictly between the thresholds,
small-decay and large-decay exponential integrands give opposite signs.
Constants are integrated exactly.
\end{proof}

\section{The refinement-new-node rule}

Define
\begin{equation}
N_{q,n}(f)
=\frac{qT_{qn}(f)-T_n(f)}{q-1}.
\label{eq:new-node}
\end{equation}
Cancellation of the coarse-grid nodes gives
\begin{equation}
N_{q,n}(f)
=
\frac{L}{n(q-1)}
\sum_{j=0}^{n-1}\sum_{r=1}^{q-1}
f\left(\frac{(j+r/q)L}{n}\right).
\label{eq:new-node-explicit}
\end{equation}
Thus \(N_{q,n}\) is the equal-weight rule supported precisely on the new
points introduced by \(q\)-fold refinement.

For \(\alpha\in\R\), let
\[
S_{n,\alpha}(f)
=\alpha T_n(f)+(1-\alpha)N_{q,n}(f),
\]
and extrapolate \(S_{n_j,\alpha}(f)\) in \(n_j^{-2}\) to zero.

\begin{corollary}[Sharp new-node blend]
\label{cor:new-node}
The extrapolated rule is a universal upper bound on \(\cC_L\) if and only if
\begin{equation}
\boxed{
\alpha\ge\frac{q^{2m-1}-1}{q^{2m}-1}.}
\label{eq:alpha-upper}
\end{equation}
It is a universal lower bound if and only if
\begin{equation}
\boxed{\alpha\le0.}
\label{eq:alpha-lower}
\end{equation}
For intermediate \(\alpha\), the error has no universal sign.
\end{corollary}

\begin{proof}
By linearity, the extrapolate equals \(Q_\beta\) with
\[
\beta=\frac{q(1-\alpha)}{q-1}.
\]
Substitution into \cref{thm:blend} gives
\eqref{eq:alpha-upper} and \eqref{eq:alpha-lower}.
\end{proof}

\begin{remark}[Midpoint and Simpson]
For \(q=2\), \(N_{2,n}\) is the composite midpoint rule \(M_n\).
Equation \eqref{eq:alpha-upper} becomes
\[
\alpha\ge
\frac{2^{2m-1}-1}{2^{2m}-1}.
\]
At \(m=1\), the lower endpoint is the midpoint rule and the upper endpoint is
Simpson's rule. For general \(m\), midpoint extrapolation is a strict lower
bound and the critical midpoint--trapezoid blend is a strict upper bound for
every nonconstant member of \(\cC_L\).
\end{remark}

\section{Scope and relation to earlier work}

The theorem concerns finite positive Laplace transforms. This class permits
unbounded derivative moments and does not require endpoint derivatives for an
Euler--Maclaurin argument. The refinement factor is an integer because the
upper constant \(q^{-1}\) comes from the spectral density of multiples of
\(q\). The mesh counts themselves are arbitrary.

Classical Romberg positivity proves one-sidedness through Peano kernels
\cite{DahlquistBjorck}. Meir--Sharma and Stroud treated the classical method
and its error estimates \cite{MeirSharma1965,Stroud1965}; Bulirsch--Stoer
developed general extrapolation bounds and extrapolatory quadrature
\cite{BulirschStoer1966,BulirschStoer1967}; and Lynch studied generalized
trapezoid formulas and Romberg errors \cite{Lynch1967}. Albrecht studied nested
interval bounds \cite{Albrecht1972}; Ström studied strict bounds and monotonicity
\cite{Strom1967,Strom1972}; Förster and Brass--Fischer developed Romberg error
bounds \cite{Foerster1982,BrassFischer1999}; Fischer established definiteness
for the Bulirsch sequence \cite{Fischer2002}; and averaged
midpoint--trapezoid inequalities have been developed in several forms,
including \cite{Liu2011}. Brass and Petras provide a modern monograph treatment
of the surrounding quadrature theory \cite{BrassPetras}. The bounded claim
here is the sharp error-ratio interval \((q^{-2m},q^{-1})\), its two independent
endpoint mechanisms, and the exact consequences in
\cref{cor:enclosure,thm:blend,cor:new-node}. This claim remains subject to
full-text review of the close archival sources.

For consecutive dyadic meshes, the upper endpoint agrees with the classical
Romberg construction. No novelty is claimed for that specialization or for
Richardson coefficients themselves.

\section{Conclusion}

The squared-mesh trapezoidal error of a finite positive Laplace transform has a
positive complete Bernstein representation. Arbitrary origin extrapolation
turns every spectral component into a positive product. Integer mesh
refinement then selects the spectral modes divisible by \(q\), placing the
refined error sharply between \(q^{-2m}\) and \(q^{-1}\) times the coarse
error. The result supplies a certified enclosure and a complete classification
of affine coarse--refined blends without geometric assumptions on the
underlying panel counts.

\begin{thebibliography}{99}

\bibitem{Albrecht1972}
J.~Albrecht,
Intervallschachtelungen beim Romberg-Verfahren,
\emph{Z. Angew. Math. Mech.} 52 (1972), 433--435.

\bibitem{BrassFischer1999}
H.~Brass and J.-W.~Fischer,
Error bounds for Romberg quadrature,
\emph{Numerische Mathematik} 82 (1999), 389--408.
\url{https://doi.org/10.1007/s002110050424}

\bibitem{BrassPetras}
H.~Brass and K.~Petras,
\emph{Quadrature Theory: The Theory of Numerical Integration on a Compact
Interval}, Mathematical Surveys and Monographs 178,
American Mathematical Society, Providence, 2011.
\url{https://doi.org/10.1090/surv/178}

\bibitem{BulirschStoer1966}
R.~Bulirsch and J.~Stoer,
Asymptotic upper and lower bounds for results of extrapolation methods,
\emph{Numerische Mathematik} 8 (1966), 93--104.
\url{https://doi.org/10.1007/BF02163179}

\bibitem{BulirschStoer1967}
R.~Bulirsch and J.~Stoer,
Numerical quadrature by extrapolation,
\emph{Numerische Mathematik} 9 (1967), 271--278.
\url{https://doi.org/10.1007/BF02162420}

\bibitem{DahlquistBjorck}
G.~Dahlquist and \AA.~Bj\"orck,
\emph{Numerical Methods in Scientific Computing, Volume I},
SIAM, Philadelphia, 2008, Chapter~5.
\url{https://jasoncantarella.com/downloads/Dahlquist-Bjorck-chapter-5.pdf}

\bibitem{Fischer2002}
J.-W.~Fischer,
Romberg quadrature using the Bulirsch sequence,
\emph{Numerische Mathematik} 90 (2002), 509--519.
\url{https://doi.org/10.1007/s002110100271}

\bibitem{Foerster1982}
K.-J.~F\"orster,
Fehlerschranken bei der Romberg-Quadratur,
\emph{Z. Angew. Math. Mech.} 62 (1982), 133--135.
\url{https://doi.org/10.1002/zamm.19820620209}

\bibitem{Liu2011}
Z.~Liu,
More on the averaged midpoint-trapezoid type rules,
\emph{Applied Mathematics and Computation} 218 (2011), 1389--1398.
\url{https://www.sciencedirect.com/science/article/pii/S0096300311008484}

\bibitem{Lynch1967}
R.~E.~Lynch,
Generalized trapezoid formulas and errors in Romberg quadrature,
in \emph{Blanch Anniversary Volume}, Aerospace Research Laboratory,
U.S. Air Force, Washington, D.C., 1967, pp.~215--229.

\bibitem{MeirSharma1965}
A.~Meir and A.~Sharma,
On the method of Romberg quadrature,
\emph{J. Soc. Indust. Appl. Math. Ser. B Numer. Anal.} 2 (1965), 250--258.
\url{https://doi.org/10.1137/0702019}

\bibitem{SchillingSongVondracek}
R.~L.~Schilling, R.~Song, and Z.~Vondra\v{c}ek,
\emph{Bernstein Functions: Theory and Applications}, 2nd ed.,
de Gruyter, Berlin, 2012.
\url{https://doi.org/10.1515/9783110269338}

\bibitem{Strom1967}
T.~Str\"om,
Strict error bounds in Romberg quadrature,
\emph{BIT} 7 (1967), 314--321.
\url{https://doi.org/10.1007/BF01939325}

\bibitem{Strom1972}
T.~Str\"om,
Monotonicity in Romberg quadrature,
\emph{Mathematics of Computation} 26 (1972), 461--465.
\url{https://www.ams.org/journals/mcom/1972-26-118/home.html}

\bibitem{Stroud1965}
A.~H.~Stroud,
Error estimates for Romberg quadrature,
\emph{J. Soc. Indust. Appl. Math. Ser. B Numer. Anal.} 2 (1965), 480--488.
\url{https://doi.org/10.1137/0702038}

\bibitem{Widder}
D.~V.~Widder,
\emph{The Laplace Transform},
Princeton University Press, Princeton, 1941.

\end{thebibliography}

\end{document}
